% !TEX root = ../main-es.tex
% Las cajas de ayuda no son contenido: texcount las ignora (ver `make wordcount`).
%TC:macro \limite [ignore,ignore]
%TC:macro \nota [ignore]
%TC:macro \TODO [ignore]
%TC:macro \parencite [ignore]
%TC:macro \textcite [ignore]
\section{Agentes}
\limite{0}{Agentes: identifique los agentes que intervienen en su proyecto. Explique
quiénes son y qué rol desempeña cada uno dentro de la dinámica del sistema.}

La dinámica del sistema la sostienen tres actores: Un \textbf{agente recolector} (el protagonista) implementado en seis arquitecturas que cubren la taxonomía de \textcite{russellNorvig2020} cap.~2, un \textbf{oráculo de replanificación} que abstrae al sistema estocástico de IA y un \textbf{entorno hostil} (los recursos móviles) que hace que el plan caduque. Existe también un cuarto papel referente al \textbf{humano autorizador} de ediciones, se modela dentro de la arquitectura gobernada con procesos de Human in The Loop (HiTL).

\begin{table}[ht]
\centering
\small
\begin{tabular}{@{}l L{4.6cm} L{5.2cm}@{}}
\toprule
\textbf{Agente} & \textbf{Quién es} & \textbf{Rol en la dinámica} \\
\midrule
Reflejo simple & Se mueve con avidez hacia el recurso visible más cercano (sin plan ni memoria) & Línea base reactiva: reevalúa en cada tick y nunca sufre un plan obsoleto (vende previsión, compra oportunismo) \\
Basado en modelo & Mantiene un mapa de creencias del mundo; reactivo, sin plan explícito & Aísla el aporte de la memoria interna frente al reflejo puro \\
Lazo abierto & Calcula una ruta A* una vez y la ejecuta sin feedback & Representa el flujo estático: ignora las divergencias entre plan y mundo \\
Reinicio en frío & Al divergir, aborta la corrida y replanifica desde cero en la siguiente & Modela el comportamiento de un sistema operativo y gestión de procesos
%  de OASys hoy (ediciones aplicadas solo en la corrida siguiente; el progreso actual se pierde)
\\
Lazo cerrado & Monitorea la ejecución y, al divergir, replanifica en sitio pagando $c$ ticks & Representa la modificación/reparación oportuna posteriori
%  el \texttt{run=hot} no implementado: la reparación oportuna
del plan con costo explícito \\
Gobernado & Como el de lazo cerrado, pero con presupuesto de $B$ ediciones y
% líneas inmutables \texttt{@lock} 
acciones inmutables
& Incorpora al humano autorizador como una restricción del sistema: Amortigua el \emph{thrashing} de reescrituras \\

\addlinespace
Oráculo & El replanificador abstracto con competencia $p$ y costo $c$ &
Reemplaza al LLM real: con probabilidad $p$ propone la reparación óptima (A* sobre el
sufijo) y con $1-p$ una reparación plausible pero equivocada \\
\bottomrule
\end{tabular}
\caption{Agentes del sistema y su papel. Las seis primeras filas son arquitecturas del
mismo agente recolector; las dos últimas completan la dinámica.}
\end{table}

La comparación entre estas políticas \emph{es} el experimento: la diferencia de
utilidad entre reinicio en frío y lazo cerrado cuantifica el costo de no haber
implementado la aplicación en caliente, y la diferencia entre lazo cerrado y gobernado
mide el valor de acotar la propia replanificación.

% \nota{Si el avance pide ``los agentes'' en sentido estricto, la tabla de seis
% arquitecturas puede resumirse a tres papeles: recolector, oráculo y entorno.}
